\documentclass[12pt]{article}
\usepackage[T1]{fontenc}
\usepackage[utf8]{inputenc}
\usepackage{lmodern}
\usepackage{textcomp}
\usepackage{enumitem}
\usepackage{geometry}
\usepackage{graphicx}
\usepackage{amsmath, amssymb}
\geometry{margin=1in}
\begin{document}
\begin{center}
Computer Science - Graduate Level Assessment 14 \
Total Marks: 40 \
Answer all questions.
\end{center}

\section*{Multiple Choice (10 marks)}
\begin{enumerate}[label=\arabic*.]
\item Which of the following is NOT a reasonable justification for choosing to busy-wait on an asynchronous event?
\begin{enumerate}[label=(\alph*)]
    \item The wait is expected to be short.
    \item A busy-wait loop is easier to code than an interrupt handler.
    \item There is no other work for the processor to do.
    \item The program executes on a time-sharing system.
\end{enumerate}
\item Which of the following sets of bit strings CANNOT be described with a regular expression?
\begin{enumerate}[label=(\alph*)]
    \item All bit strings whose number of zeros is a multiple of five
    \item All bit strings starting with a zero and ending with a one
    \item All bit strings with an even number of zeros
    \item All bit strings with more ones than zeros
\end{enumerate}
\item Which of the following are most vulnerable to injection attacks?
\begin{enumerate}[label=(\alph*)]
    \item Session IDs
    \item Registry keys
    \item Network communications
    \item SQL queries based on user input
\end{enumerate}
\item You are given a message (m) and its OTP encryption (c). Can you compute the OTP key from m and c ?
\begin{enumerate}[label=(\alph*)]
    \item No, I cannot compute the key.
    \item Yes, the key is k = m xor c.
    \item I can only compute half the bits of the key.
    \item Yes, the key is k = m xor m.
\end{enumerate}
\item A \_\_\_\_\_\_\_\_\_\_\_ is a method in which a computer security mechanism is bypassed untraceable for accessing the computer or its information.
\begin{enumerate}[label=(\alph*)]
    \item front-door
    \item backdoor
    \item clickjacking
    \item key-logging
\end{enumerate}
\end{enumerate}

\section*{True / False (10 marks)}
\textit{Answer True or False}
\begin{enumerate}[label=\arabic*.]
\setcounter{enumi}{5}
\item True or False: Mutation-based fuzzing exclusively applies to file-based inputs and is not suitable for network protocol testing.
\item True or False: The decimal number 0.1 can be represented exactly in binary notation without approximation.
\item True or False: WPA 2 is a recent protocol that simplifies the process of connecting devices to wireless networks.
\item True or False: Encryption and decryption provide both confidentiality and integrity by ensuring that data is secure and unaltered.
\item True or False: A packet filter firewall functions at the application or transport layer, analyzing data payloads for filtering.
\end{enumerate}

\section*{Long Form Response (20 marks)}
\textit{Answer each question with a clear and complete explanation.}
\begin{enumerate}[label=\arabic*.]
\setcounter{enumi}{10}
\item Write a python function to check whether the given number is a perfect square or not.
\item Write a python function to find the minimum operations required to make two numbers equal.
\end{enumerate}

\vfill
\noindent\textit{End of Paper}
\end{document}